\section{Proof status and verification}\label{sec:verification}

The Lean-verified results are exact homogenization
(\cref{prop:inflation}) and the dimension interfaces
(\cref{lem:kernel-range}). The shared
\leanfile{LeanMachine/HomogeneousInflation.lean} module constructs
the two checks and their coordinate permutation, proves the kernel
equivalence, and applies Mathlib's finite-product dimension formula.
\leanfile{LeanMachine/KernelRangeBound.lean} reuses Mathlib's
dimension monotonicity, rank-nullity, quotient dimension, and span
closure for the dimension, quotient-deficit, and split-generator bounds.
The paper modules expose these results without duplicating their proofs.
The English outlines are guides; the Lean declarations are authoritative.

The remaining graph, finite and asymptotic cover, nonlinear-source, and group-word
results have complete written proofs.
In particular, the local-generation theorem is not described as
formalized merely because its final dimension inference is checked.
The native hypotheses, free-group separation, imported Cohn and
Jaikin-Zapirain--Souza theorems, parameter uniformization, coefficient
filtrations, weighted determinant argument, and full-matrix-algebra
inflation example are not formalized in this manuscript. Nor is the
written product-subcode corollary part of the formal kernel equivalence.
The source-coset argument and central-layer group-word specialization are
also written proofs, not additional claims about the scope of that interface.
The README's formalization todos record the remaining dependencies and
existing library material to reuse. No floating-point estimate or
experimental observation is used as a proof premise.

The finite Clebsch lift, reciprocal-block, and small-representation
arguments are complete written proofs. Their algebraic base signing
rigidity proof does not depend on a finite rank computation. The separate
48-vertex seed uses the exact witness in
\path{tools/clebsch-seed-certificate.json}: its literal kernel basis,
original signing-map minor, and explicit inverse are verified by
\texttt{python3 tools/check\_clebsch\_seed.py}. This finite identity
certificate is part of the written proof and is not described as
kernel-checked Lean code. It uses no floating-point arithmetic or solver.

Lean and Mathlib are pinned to v4.34.1. The paper requires the
sibling shared package through a relative path. Clone the whole
repository with its directories intact. From the paper directory,
\texttt{python3 tools/lean\_index.py} generates the declaration
index and \texttt{lean/PaperAxioms.lean} from the same annotations.
In \texttt{lean/}, run \texttt{lake exe cache get},
\texttt{lake build}, and \texttt{lake env lean PaperAxioms.lean}.
The cited declarations use at most \texttt{propext},
\texttt{Classical.choice}, and \texttt{Quot.sound}.
In \texttt{paper/}, run \texttt{latexmk -pdf main.tex}.
